% Generated draft: requires mathematical and source review.
\documentclass[11pt]{article}
\usepackage[margin=25mm]{geometry}
\usepackage{fontspec}
\setmainfont{Latin Modern Roman}
\usepackage{amsmath,amssymb,mathtools,unicode-math}
\setmathfont{Latin Modern Math}
\usepackage{enumitem,needspace,xurl,hyperref}
\setlength{\parindent}{0pt}
\setlength{\parskip}{5pt}
\setlength{\emergencystretch}{3em}
\newcommand{\N}{\mathbb{N}}
\newcommand{\Z}{\mathbb{Z}}
\newcommand{\Q}{\mathbb{Q}}
\newcommand{\R}{\mathbb{R}}
\newcommand{\C}{\mathbb{C}}
\newcommand{\F}{\mathbb{F}}
\newcommand{\A}{\mathbb{A}}
\newcommand{\PP}{\mathbb{P}}
\newcommand{\E}{\mathbb{E}}
\newcommand{\eps}{\varepsilon}
\newcommand{\dd}{\,\mathrm{d}}
\begin{document}
\section*{UnsolvedMath ID: 701666}
\textbf{Source problem code:} catalogue entry
\subsection{Definitions and mathematical statement}
A free group has finite rank if it has a finite free basis; its rank is the cardinality of such a basis. For a possibly infinitely generated free subgroup \(L\), let \(\operatorname{rank}(L)\) denote the cardinality of a free basis. In particular, an infinite rank is greater than every finite rank.

Let \(K\leq H\leq F\) be free groups of finite rank satisfying
\[
\operatorname{rank}(H)<\operatorname{rank}(K)
\]
and
\[
\forall L\leq H,\qquad
\bigl(K\leq L\ \text{and}\ L\ne H\bigr)
\Longrightarrow
\operatorname{rank}(L)\geq\operatorname{rank}(K).
\]
The quantifier over \(L\) includes all intermediate subgroups, whether or not they are finitely generated, and allows \(L=K\).

Let \(\iota:H\hookrightarrow F\) be the inclusion homomorphism. Must
\[
\forall\varphi\in\operatorname{Hom}(H,F),\qquad
\left(\forall k\in K,\ \varphi(k)=k\right)
\Longrightarrow \varphi=\iota
\]
hold? Here \(\varphi\) is an arbitrary group homomorphism, with no injectivity or surjectivity assumption.

The source writes \(K<H<F\) without defining the symbol \(<\). If it denotes proper subgroup inclusion, impose the additional hypothesis \(H\ne F\); the rank inequality already implies \(K\ne H\). If it denotes subgroup inclusion allowing equality, no additional hypothesis is imposed. This convention remains to be checked against the cited source.
\subsection{Short English statement}
Suppose a finite-rank free subgroup \(K\) lies in a free group \(H\) of smaller finite rank, and every proper subgroup of \(H\) containing \(K\) has rank at least that of \(K\). For a finite-rank free ambient group \(F\) containing \(H\), must every homomorphism \(H\to F\) that fixes \(K\) pointwise be the inclusion? Whether the source requires \(H\) to be proper in \(F\) remains a notation qualification.
\subsection{Sources}
\begin{itemize}
\item Local UnsolvedMath registry, supplied complete source record: permanent ID \(701666\), exact source problem code \texttt{catalogue entry}, title \emph{Kourovka Notebook Problem 20.8}.
\item The supplied record attributes the problem to G. M. Bergman, with proposed year 2022, and cites \emph{The Kourovka Notebook (Unsolved Problems in Group Theory)}, 20th issue (2022), original problem number 20.8, and the 21st issue edition (2026), \texttt{arXiv:1401.0300v48}. Its collection description names E. I. Khukhro and V. D. Mazurov as editors.
\item The registry labels the problem \emph{open}. This is supplied source metadata, not an independently verified mathematical or literature assessment. The cited editions and attribution have not been independently checked for this draft.
\end{itemize}
\end{document}
